% !TEX program = xelatex
% 琼台师范学院本科毕业论文（设计）LaTeX 模板
% 说明：学校 2026 届材料需验证码下载，本模板按《琼台师范学院本科生毕业论文（设计）
% 工作规定（2021 年修订版）》的装订顺序和通用本科论文格式搭建。
% 使用前请务必以学院当年下发的 Word 模板核对页边距、字号和封面字段。
\documentclass[12pt,a4paper,fontset=none]{ctexart}

%====================== 页面与字体 ======================
\usepackage{geometry}
\geometry{top=25.4mm,bottom=25.4mm,left=30mm,right=25mm}
\usepackage{fontspec}
\usepackage{xeCJK}
% MIT; default preview uses TeX-distributed fonts, NOT school typography.
\RequirePackage{fontspec,xeCJK,etoolbox}
\providecommand{\ThesisFontMode}{preview}
\ifdefstring{\ThesisFontMode}{school}{%
  \newcommand{\ThesisCheckFont}[1]{\IfFontExistsTF{#1}{}{\PackageError{thesis-fonts}{Missing font: #1}{Install legally licensed fonts or use THESIS_FONTS=preview.}}}
  \ThesisCheckFont{Times New Roman}\ThesisCheckFont{SimSun}
  \ThesisCheckFont{SimHei}\ThesisCheckFont{FangSong}\ThesisCheckFont{KaiTi}
  \setmainfont{Times New Roman}
  \setCJKmainfont[AutoFakeBold=2]{SimSun}
  \setCJKsansfont[AutoFakeBold=2]{SimHei}
  \setCJKmonofont{FangSong}
  \setCJKfamilyfont{zhsong}[AutoFakeBold=2]{SimSun}
  \setCJKfamilyfont{zhhei}[AutoFakeBold=2]{SimHei}
  \setCJKfamilyfont{zhkai}[AutoFakeBold=2]{KaiTi}
  \setCJKfamilyfont{zhfs}{FangSong}
}{%
  \ifdefstring{\ThesisFontMode}{preview}{}{\PackageError{thesis-fonts}{Invalid font mode}{Use preview or school.}}
  \setmainfont{texgyretermes-regular.otf}[BoldFont=texgyretermes-bold.otf,ItalicFont=texgyretermes-italic.otf,BoldItalicFont=texgyretermes-bolditalic.otf]
  \setCJKmainfont{FandolSong-Regular.otf}[BoldFont=FandolSong-Bold.otf]
  \setCJKsansfont{FandolHei-Regular.otf}[BoldFont=FandolHei-Bold.otf]
  \setCJKmonofont{FandolFang-Regular.otf}
  \setCJKfamilyfont{zhsong}{FandolSong-Regular.otf}[BoldFont=FandolSong-Bold.otf]
  \setCJKfamilyfont{zhhei}{FandolHei-Regular.otf}[BoldFont=FandolHei-Bold.otf]
  \setCJKfamilyfont{zhkai}{FandolKai-Regular.otf}[AutoFakeBold=2]
  \setCJKfamilyfont{zhfs}{FandolFang-Regular.otf}
}
% Geometric/shapes may not exist in the serif font; ctex's open symbol font.
\xeCJKDeclareCharClass{CJK}{"2460 -> "24FF,"3200 -> "32FF}

\newcommand{\song}{\CJKfamily{zhsong}}
\newcommand{\hei}{\CJKfamily{zhhei}}
\newcommand{\kai}{\CJKfamily{zhkai}}
\newcommand{\fs}{\CJKfamily{zhfs}}

\usepackage{amsmath,amssymb}
\usepackage{graphicx}
\usepackage{booktabs}
\usepackage{caption}
\usepackage{titlesec}
\usepackage{tocloft}
\usepackage{fancyhdr}
\usepackage{indentfirst}
\usepackage[hidelinks]{hyperref}
\usepackage{gbt7714}
\usepackage{tabularx,longtable,listings,perpage,tikz,pgfplots}
\usetikzlibrary{positioning,arrows.meta}
\pgfplotsset{compat=1.18}
\xeCJKDeclareCharClass{CJK}{"2460 -> "24FF}
\MakePerPage{footnote}
\renewcommand{\thefootnote}{\ifcase\value{footnote}\or①\or②\or③\or④\or⑤\or⑥\or⑦\or⑧\or⑨\or⑩\or⑪\or⑫\or⑬\or⑭\or⑮\or⑯\or⑰\or⑱\or⑲\or⑳\else\arabic{footnote}\fi}

%====================== 论文信息（改这里） ======================
\newcommand{\thesistitle}{小学语文课堂提问策略研究}
\newcommand{\thesistitleeng}{A Study on Questioning Strategies in Primary School Chinese Classes}
\newcommand{\thesisauthor}{李\quad 四}
\newcommand{\thesisid}{2022010101}
\newcommand{\thesisschool}{文学院}
\newcommand{\thesismajor}{汉语言文学}
\newcommand{\thesisclass}{2022 级 1 班}
\newcommand{\thesisadvisor}{王\quad 五}
\newcommand{\thesisdate}{二〇二六年五月}

%====================== 标题 ======================
\setcounter{secnumdepth}{3}
\setcounter{tocdepth}{3}
\renewcommand{\thesection}{\arabic{section}}
\renewcommand{\thesubsection}{\thesection.\arabic{subsection}}
\renewcommand{\thesubsubsection}{\thesubsection.\arabic{subsubsection}}

\titleformat{\section}{\hei\zihao{3}}{\thesection}{0.6em}{}
\titleformat{\subsection}{\hei\zihao{4}}{\thesubsection}{0.6em}{}
\titleformat{\subsubsection}{\hei\zihao{-4}}{\thesubsubsection}{0.6em}{}
\titlespacing*{\section}{0pt}{0.6em}{0.4em}
\titlespacing*{\subsection}{0pt}{0.4em}{0.3em}
\titlespacing*{\subsubsection}{0pt}{0.3em}{0.2em}

%====================== 目录 ======================
\renewcommand{\contentsname}{目\quad 录}
\renewcommand{\cfttoctitlefont}{\hfill\hei\zihao{3}}
\renewcommand{\cftaftertoctitle}{\hfill\mbox{}}
\renewcommand{\cftsecfont}{\song\zihao{-4}}
\renewcommand{\cftsecpagefont}{\song\zihao{-4}}
\renewcommand{\cftsubsecfont}{\song\zihao{-4}}
\renewcommand{\cftsubsecpagefont}{\song\zihao{-4}}
\renewcommand{\cftsubsubsecfont}{\song\zihao{-4}}
\renewcommand{\cftsubsubsecpagefont}{\song\zihao{-4}}
\renewcommand{\cftdotsep}{\cftnodots}
\setlength{\cftbeforesecskip}{0.2em}
\setlength{\cftbeforesubsecskip}{0pt}
\setlength{\cftbeforesubsubsecskip}{0pt}
\setlength{\cftsecindent}{0pt}
\setlength{\cftsubsecindent}{1.5em}
\setlength{\cftsubsubsecindent}{3em}

%====================== 页眉页脚 ======================
\fancypagestyle{main}{%
  \fancyhf{}%
  \fancyfoot[C]{\song\zihao{5}\thepage}%
  \renewcommand{\headrulewidth}{0pt}%
}
\pagestyle{main}

%====================== 图表公式 ======================
\DeclareCaptionFont{wu}{\song\zihao{5}}
\captionsetup{font=wu,labelsep=space,justification=centering}
\captionsetup[table]{position=top}
\renewcommand{\thefigure}{\arabic{section}.\arabic{figure}}
\renewcommand{\thetable}{\arabic{section}.\arabic{table}}

% 可换行封面字段；具体封面版式仍待当届官方原件校正。
\newcommand{\coverfield}[1]{\parbox[c]{9.6cm}{\centering #1\par\kern2pt\hrule height0.4pt}}
\counterwithin{figure}{section}
\counterwithin{table}{section}
\counterwithin{equation}{section}
\begin{document}
\linespread{1.5}\zihao{-4}\selectfont
\setlength{\parindent}{2em}
\setlength{\parskip}{0pt}

%====================== 封面 ======================
\begin{titlepage}
\thispagestyle{empty}
\begin{center}
\vspace*{1.0cm}
{\hei\zihao{1}\bfseries 琼台师范学院}\\[1.6cm]
{\hei\zihao{2} 本科毕业论文（设计）}\\[2.2cm]
\renewcommand{\arraystretch}{2.0}
\begin{tabular}{r@{\hspace{0.3em}}l}
{\hei\zihao{4} 题\quad 目：} & \coverfield{\hei\zihao{3}\thesistitle} \\[0.2cm]
{\hei\zihao{4} 姓\quad 名：} & \coverfield{\kai\zihao{-3}\thesisauthor} \\[0.2cm]
{\hei\zihao{4} 学\quad 号：} & \coverfield{\kai\zihao{-3}\thesisid} \\[0.2cm]
{\hei\zihao{4} 学\quad 院：} & \coverfield{\kai\zihao{-3}\thesisschool} \\[0.2cm]
{\hei\zihao{4} 专\quad 业：} & \coverfield{\kai\zihao{-3}\thesismajor\ \thesisclass} \\[0.2cm]
{\hei\zihao{4} 指导教师：} & \coverfield{\kai\zihao{-3}\thesisadvisor} \\
\end{tabular}
\end{center}

\vfill
\begin{center}
{\hei\zihao{3}\thesisdate}
\end{center}
\end{titlepage}

%====================== 独创性声明 ======================
\newpage
\thispagestyle{empty}
\begin{center}
{\hei\zihao{2} 独创性声明（占位）}
\end{center}
\vspace{0.8em}
{\song\zihao{-4}\hspace*{2em}
本页仅预留独创性声明位置，尚未取得当届学校声明原文。正式提交时须自行从学院取得并替换，不应将本段提示或通用声明误认为学校审定文字。作者应确认引用、成果与贡献真实准确，并按学院要求本人签署。}

\vspace{2.5cm}
\begin{flushright}
作者签名：\underline{\makebox[3cm]{}}\\[1.2em]
日\quad 期：\underline{\makebox[3cm]{}}\ 年\ \underline{\makebox[1cm]{}}\ 月\ \underline{\makebox[1cm]{}}\ 日
\end{flushright}

%====================== 授权声明 ======================
\newpage
\thispagestyle{empty}
\begin{center}
{\hei\zihao{2} 授权声明（占位）}
\end{center}
\vspace{0.8em}
{\song\zihao{-4}\hspace*{2em}
本页仅预留论文使用授权声明位置，不产生任何实际授权。正式提交前请取得当届原文，理解授权对象、范围、保密安排及本人权利后按要求填写。本项目不代替作者同意，也不捏造学校授权条款。}

\vspace{2.5cm}
\begin{flushright}
作者签名：\underline{\makebox[3cm]{}}\qquad 指导教师签名：\underline{\makebox[3cm]{}}\\[1.2em]
日\quad 期：\underline{\makebox[3cm]{}}\ 年\ \underline{\makebox[1cm]{}}\ 月\ \underline{\makebox[1cm]{}}\ 日
\end{flushright}

%====================== 中文摘要 ======================
\newpage
\pagenumbering{roman}
\pagestyle{main}
\addcontentsline{toc}{section}{摘要}
\begin{center}
{\hei\zihao{-2}\bfseries \thesistitle}
\end{center}
\vspace{0.6em}
\begin{center}
{\hei\zihao{3} 摘\quad 要}
\end{center}
\vspace{0.4em}

{\song\zihao{-4}
本文件为教育研究论文写作与排版示例，不代表已经开展课堂观察或访谈。模板展示研究问题、研究设计、资料编码、结果报告与讨论的组织方法，提供观察维度表、流程图、公式、脚注、附录等可编辑元素。没有实际收集三所学校、二十四节课或十二位教师的数据，也不提供虚构的伦理批准或统计结论。正式论文应使用真实且获准使用的资料，明确样本选择、知情同意、隐私保护和分析过程。当前只核实了琼台师范学院公开工作规定中的装订结构，尚未取得当年官方撰写规范，因此封面字段、页边距、字号及页码样式仍须与学院要求核对。}

\vspace{0.6em}
\noindent{\hei\zihao{-4} 关键词：}{\song\zihao{-4}小学语文；课堂提问；阅读教学；教学策略}

%====================== 英文摘要 ======================
\newpage
\addcontentsline{toc}{section}{Abstract}
\begin{center}
{\rmfamily\zihao{-2}\bfseries \thesistitleeng}
\end{center}
\vspace{0.6em}
\begin{center}
{\rmfamily\zihao{3}\bfseries Abstract}
\end{center}
\vspace{0.4em}

{\rmfamily\zihao{-4}
This document demonstrates the structure and typesetting of an educational research thesis. It does not report classroom observations, teacher interviews, or empirical findings. The examples cover research questions, sampling, coding, tables, diagrams, equations, and supplementary material. Replace all examples with independently conducted and appropriately approved research. The layout remains provisional because the current official writing specification has not been obtained.}

\vspace{0.6em}
\noindent{\rmfamily\zihao{-4}\bfseries Key words:} {\rmfamily\zihao{-4}primary school Chinese; classroom questioning; reading instruction; teaching strategy}

%====================== 目录 ======================
\newpage
\tableofcontents
\thispagestyle{empty}

%====================== 正文 ======================
\newpage
\pagenumbering{arabic}
\pagestyle{main}

\section{绪论}
\subsection{问题来源与研究目标}
教育研究应把一般教学关切转化为具体、可回答的问题。例如“如何识别课堂提问的层次与反馈方式”比“如何提高教育质量”具有更明确的观察边界。选题时还需确认研究进入现场的条件、时间安排与资料使用权限。
\subsection{概念界定与文献综述}
将课堂提问、候答时间、追问及反馈等概念作操作性定义。文献综述应比较已有研究的方法与证据，不能将理论分类直接当成效果结论。本例引用真实的教育分类和质性分析著作，演示来源标注\cite{anderson2001,braun2006}；它们不构成本选题的完整综述。
\subsubsection{研究边界}
明确学段、学科、课型与分析单位，说明哪些问题不在研究范围内。课堂记录包含未成年人信息时，应按机构要求取得必要审批、监护人同意和参与者适当同意，不应凭模板文字声称已经获批\footnote{本脚注演示补充说明，模板没有收集任何学生或教师资料。}。
\section{研究设计}
\subsection{抽样与资料来源}
说明学校、教师和课堂的选择方式以及纳入、排除标准。便利抽样不能自动代表总体，不同课堂的观察次数也不等于独立教师数量。正文只记录实际完成的工作，不预填访谈人数、课时数量或场所名称。
\subsection{研究流程}
图~\ref{fig:education} 给出一条可供修改的工作路线，不是学校要求的固定方法。
\begin{figure}[htbp]\centering
\begin{tikzpicture}[node distance=6mm,b/.style={draw,minimum width=65mm,minimum height=9mm,font=\small},>=Stealth]
\node[b](a){界定问题与取得必要许可};
\node[b,below=of a](b){试用观察工具并修订编码规则};
\node[b,below=of b](c){收集资料与匿名化整理};
\node[b,below=of c](d){分析、复核与报告局限};
\draw[->](a)--(b);\draw[->](b)--(c);\draw[->](c)--(d);
\end{tikzpicture}
\caption{研究组织流程示例}\label{fig:education}
\end{figure}
\subsection{观察维度与编码规则}
表~\ref{tab:coding} 只展示编码表结构。正式研究需通过试编码明确类别边界，处理难以判断的情况，并记录规则修改。
\begin{table}[htbp]\centering\caption{观察维度示例（非调查结果）}\label{tab:coding}
\begin{tabularx}{\linewidth}{lXX}\toprule 维度 & 记录内容 & 注意事项\\\midrule
问题类型 & 问题原文与分类依据 & 不以关键词机械判断认知水平\\
候答时间 & 提问结束至首次回应的间隔 & 明确起止点与计时规则\\
教师反馈 & 追问、转问、解释等行为 & 保留上下文，允许多重编码\\
参与方式 & 自愿、点名或集体回答 & 不公开学生身份\\\bottomrule\end{tabularx}
\end{table}
\subsection{分析方法与一致性}
若使用分类统计，应说明分析单位、分母和缺失处理。例如某类问题占比可写为：
\begin{equation}\label{eq:proportion}p_k=\frac{n_k}{N}\times100\%.\end{equation}
式中，$n_k$ 是第 $k$ 类事件数，$N$ 是满足统计条件的总事件数。事件数不等于独立样本量，重复课堂不能任意当作独立观测。质性分析应呈现编码形成过程和反例，不能只用少量摘录印证既定观点。
\section{结果报告的组织方式}
\subsection{描述性结果}
实际研究应先报告资料规模、缺失和排除情况，再呈现与问题对应的结果。本模板不提供真实数据，因此没有柱状图、比例或访谈结论可报告。图表应有正文引导，资料摘录应匿名并注明内部标识。
\subsection{证据与解释的区分}
分开呈现观察事实、受访者的观点和研究者的解释。单次课堂中的沉默可能有多种原因，不能直接等同于学生不会思考；受访者的自述也不等于课堂行为的客观记录。
\section{讨论与建议}
\subsection{与已有研究对话}
比较研究发现与文献中的一致和不一致之处，解释方法差异是否可能影响结论。建议应由证据支持，并说明资源、教师经验及课型等实施条件，避免把一般常识包装成独立研究发现。
\subsection{局限与研究伦理}
说明抽样范围、观察者效应、编码主观性及资料不完整等限制。保存录音、视频和原始记录时应控制访问权限，不能把个人信息随论文公开。最终提交前还应核对学院规定的论文格式与归档材料。

\section{结论}
本模板展示研究论文的组织方式和排版元素，没有开展实证研究，不能据此判断课堂提问的实际效果。正式结论应回答研究问题，说明证据支持的发现及其适用范围，不将个别观察外推为普遍规律。

%====================== 参考文献 ======================
\newpage
\begingroup
\renewcommand{\bibsection}{%
  \phantomsection
  \addcontentsline{toc}{section}{参考文献}%
  \begin{center}\hei\zihao{3}参考文献\end{center}
  \vspace{0.5\baselineskip}%
}
\renewcommand{\bibfont}{\song\zihao{5}}
\bibliographystyle{gbt7714-2015-numeric}
\bibliography{refs}
\endgroup

%====================== 附录 ======================
\newpage
\section*{附\quad 录}
\addcontentsline{toc}{section}{附录}
{\song\zihao{5}\linespread{1.0}\selectfont
附录 A\quad 课堂观察记录表\par
附录 B\quad 教师访谈提纲\par
\par\medskip
\begin{longtable}{lp{.65\linewidth}}\toprule 字段 & 填写说明\\\midrule\endhead
课次标识 & 匿名编号，不使用教师或学生姓名\\
时间点 & 记录事件发生的相对时间\\
问题原文 & 在取得适当许可的前提下摘录\\
编码与依据 & 分类结果、上下文与疑难说明\\\bottomrule\end{longtable}
\noindent 访谈提纲示例：\par
1. 您如何为本节课设计核心问题？\par
2. 您怎样判断何时追问、转问或延长候答？\par
3. 是否有与原先预期不一致的课堂情况？\par
这些问题仅供工具设计参考；使用前须经过试访与必要的伦理审核，不是已经实施的调查记录。

}

%====================== 致谢 ======================
\newpage
\section*{致\quad 谢}
\addcontentsline{toc}{section}{致谢}
感谢指导教师在选题、调研和论文写作过程中给予的悉心指导，感谢参与观察与访谈的各位教师和同学，感谢家人一直以来的支持。

\vspace{2cm}
\begin{flushright}
作者签名：\underline{\makebox[3cm]{}}\\[1.2em]
日\quad 期：\underline{\makebox[3cm]{}}\ 年\ \underline{\makebox[1cm]{}}\ 月\ \underline{\makebox[1cm]{}}\ 日
\end{flushright}

\end{document}
